\documentclass[a4paper]{article}

\usepackage{graphicx}
\usepackage{amsmath,amssymb}
\usepackage{minted}
\usepackage{multirow}
\usepackage{float}
\usepackage{todonotes}
\usepackage{forest}
\usepackage{xstring}
\usepackage{xcolor}
\usepackage[most]{tcolorbox}
\usepackage{xparse}
\usepackage{natbib}

\usetikzlibrary{graphs,graphdrawing}
\usegdlibrary{layered}

% for external graphviz
\input{/home/donkarlo/Dropbox/repo/nd_ai_project/src/nd_ai/knoweldge_representation/language/formal/computer/programming/tex/latex/code_clips/external_graphviz.tex}

% for tags
\input{/home/donkarlo/Dropbox/repo/nd_ai_project/src/nd_ai/knoweldge_representation/language/formal/computer/programming/tex/latex/parts/control_sequence/kind/semtag/semtag.tex}

% for links
\input{/home/donkarlo/Dropbox/repo/nd_ai_project/src/nd_ai/knoweldge_representation/language/formal/computer/programming/tex/latex/code_clips/seehere.tex}

\title{Interactional linguistic}
\date{}
\author{Mohammad Rahmani}

\begin{document}
\maketitle
    Interactional linguistics (IL) is an interdisciplinary approach to grammar and interaction in the field of linguistics, that applies the methods of Conversation Analysis to the study of linguistic structures, including syntax, phonetics, morphology, and so on. Interactional linguistics is based on the principle that linguistic structures and uses are formed through interaction and it aims at understanding how languages are shaped through interaction. The approach focuses on temporality, activity implication and embodiment in interaction \seeherefooter{https://en.wikipedia.org/wiki/Interactional_linguistics}.
%    \bibliographystyle{plainnat}
%    \bibliography{/home/donkarlo/Dropbox/repo/nd_shared_research/bibliography/bibliography.bib}
\end{document}